\chapter{Conclusiones}
\thispagestyle{fancy}
\label{ch:conclusiones}



Para concluir este análisis sobre la densidadesde los sólidos medidos, se 
plasma que:

\section{Eficacia de la Estrategia Experimental}

Se afirma que la precisión de las magnitudes indirectas se relacionan fuertemente
con la geometría del objeto, además El uso del micrómetro para el 
espesor de la moneda, cuya apreciación es de $0,001 \text{ cm}$, 
permitió que este objeto presentara uno de los errores relativos 
en volumen más bajos ($1,02\%$), a pesar de las irregularidades 
superficiales inherentes al acuñado. Esto demuestra que la técnica
 de variar los puntos de contacto y controlar la presión del tornillo 
 (sección \ref{ch:montaje}) 
fue efectiva para mitigar errores aleatorios.

\section{Vulnerabilidades}

El método de propagación de errores utilizado señala que el cilindor hueco 
es el sólido más complejo de medir, incluso más que el corcho, que es un 
objeto que posee ciertas características irregulares en su cuerpo. Esto se 
debe a que se requiere de más mediciones hechas en el diámetro interno, y a 
su vez, esto requiere de una técnica más experimentada según los resultados obtenidos.
Los resultados sugieren que medir este diámetro interno requiere de más 
habilidad que el pesaje del corcho, puesto que instrumentalmente, el vernier
posee mayor apreciación que la balanza utilizada; destacando que esta última 
no estaba en perfecto estado y pudo haber interferido en lograr mejores 
resultados. 

El hallazgo más crítico de esta investigación es la densidad 
obtenida para el cilindro dorado ($\rho = 27,76 \pm 0.34 \text{ g/cm}^3$). 
Dado que este valor excede la densidad del Osmio ($\approx 22,6 \text{ g/cm}^3$), 
se concluye la existencia de un error sistemático severo originado en la balanza 
granataria o en la técnica de pesaje para masas superiores a los 
$40 \text{ g}$. Asimismo, el método de pesaje por diferencia aplicado 
al corcho, aunque ingenioso para evitar la baja sensibilidad de la balanza, 
no fue suficiente para compensar el sesgo de la misma, 
resultando en una densidad superior al estándar bibliográfico.

En términos generales se cumplió la carectización física. Sin embargo, no se 
corrigió un sesgo sistemático en el montaje inicial, dejando errores relativos mayores
al $1,00\%$. Debido a esto, se recomienda en futuras prácticas una mejor 
atención al estado de la balanza y una mejor inclusión del micrométro en las dimensiones
de los diámetros para reducir la incertidumbre dentro del sistema. 